\textbf{- Reference:} 
\href{https://mysite.science.uottawa.ca/mnevins/papers/StephenHarrigan2017LWE.pdf}{Lattice-Based Cryptography and the Learning with
Errors Problem}~\cite{lattice-crypto}


$ $

Lattice-based cryptography provides candidates for post-quantum security, under precisely specified computational assumptions and parameter choices. Resistance to every quantum attack is not proved unconditionally. This section introduces the LWE (Learning with Errors), RLWE (Ring Learning with Errors), GLWE, GLev, and GGSW constructions used later.





\subsection{Overview}
\label{subsec:lattice-overview}

All equations in this overview are over $\mathbb Z_q$. Masks are sampled uniformly and independently; errors are small integers from a specified distribution, reduced modulo $q$ only for arithmetic.

Suppose we have a single unknown $k$-dimensional vector $\vec{s}$ as a secret key, many publicly known $k$-dimensional vectors $\vec{a}^{\langle i \rangle}$.

And suppose we have a large set of the following dot products $\vec{s} \cdot \vec{a}^{\langle i \rangle}$: 

$\vec{s} \cdot \vec{a}^{\langle 0 \rangle} = s_0\cdot a_{0}^{\langle 0 \rangle} + s_1\cdot a_{1}^{\langle 0 \rangle} + \cdots + s_{k-1}\cdot a_{k-1}^{\langle 0 \rangle} = b^{\langle 0 \rangle}$

$\vec{s} \cdot \vec{a}^{\langle 1 \rangle} = s_0\cdot a_{0}^{\langle 1 \rangle} + s_1\cdot a_{1}^{\langle 1 \rangle} + \cdots + s_{k-1}\cdot a_{k-1}^{\langle 1 \rangle} = b^{\langle 1 \rangle}$

$\vec{s} \cdot \vec{a}^{\langle 2 \rangle} = s_0\cdot a_{0}^{\langle 2 \rangle} + s_1\cdot a_{1}^{\langle 2 \rangle} + \cdots + s_{k-1}\cdot a_{k-1}^{\langle 2 \rangle} = b^{\langle 2 \rangle}$

\text{ } $\vdots$

$ $

In the noiseless case, $k$ samples determine $\vec s$ if their coefficient matrix is invertible over $\mathbb Z_q$. For prime $q$, this means full rank; over a composite modulus the determinant must be a unit. Merely having $k$ equations does not ensure independence or unique recovery.

However, suppose that in each equation above, we randomly add an unknown small noise $e^{\langle i \rangle}$ (i.e., error) as follows: 

$b^{\langle 0 \rangle} = s_0\cdot a_{0}^{\langle 0 \rangle} + s_1\cdot a_{1}^{\langle 0 \rangle} + \cdots + s_{k-1}\cdot a_{k-1}^{\langle 0 \rangle} + e^{\langle 0 \rangle} \pmod q$

$b^{\langle 1 \rangle} = s_0\cdot a_{0}^{\langle 1 \rangle} + s_1\cdot a_{1}^{\langle 1 \rangle} + \cdots + s_{k-1}\cdot a_{k-1}^{\langle 1 \rangle} + e^{\langle 1 \rangle} \pmod q$

$b^{\langle 2 \rangle} = s_0\cdot a_{0}^{\langle 2 \rangle} + s_1\cdot a_{1}^{\langle 2 \rangle} + \cdots + s_{k-1}\cdot a_{k-1}^{\langle 2 \rangle} + e^{\langle 2 \rangle} \pmod q$

\text{ } $\vdots$


$ $

Noise prevents simply treating these as exact linear equations for $\vec s$. If each of $k$ errors has $r$ possible values, enumerating all $r^k$ error vectors is one possible attack. This is an upper bound for that particular enumeration, not a lower bound for all algorithms and not a security proof. Additional samples may help an attacker and can always be ignored. LWE hardness is a computational assumption for specified dimension, modulus, secret and error distributions, and sample count; sufficiently small or structured errors can make some instances easy.

Based on this intuition, the mathematical hard problem that constitutes lattice-based cryptography is as follows:

$ $

\begin{tcolorbox}[title={\textbf{\tboxlabel{\ref*{subsec:lattice-overview}} The LWE (Learning with Errors) and RLWE Problems}}]
\textbf{\underline{LWE Problem}}

Consider samples of the form: $b = \vec{s} \cdot\vec{a} + e$ (where $e$ is a small noise to be explained later). 

For each sample, choose $\vec a$ uniformly from $\mathbb Z_q^k$ and independently draw a small integer error $e$ from an explicitly specified distribution $\chi$, then reduce $e$ modulo $q$. The secret $\vec s$ is sampled once from its specified distribution and reused. The error is not uniform over $\mathbb Z_q$. A security experiment permits a polynomial number of such samples in the security parameter.

$(\vec{a}^{\langle 1 \rangle}, b^{\langle 1 \rangle})$, where $b^{\langle 1 \rangle} = \vec{s} \cdot \vec{a}^{\langle 1 \rangle} + e^{\langle 1 \rangle}$ 

$(\vec{a}^{\langle 2 \rangle}, b^{\langle 2 \rangle})$, where $b^{\langle 2 \rangle} = \vec{s} \cdot \vec{a}^{\langle 2 \rangle} + e^{\langle 2 \rangle}$ 

$(\vec{a}^{\langle 3 \rangle}, b^{\langle 3 \rangle})$, where $b^{\langle 3 \rangle} = \vec{s} \cdot \vec{a}^{\langle 3 \rangle} + e^{\langle 3 \rangle}$ 

\text{ } $\vdots$

Suppose that the attacker has a sufficiently large number of $(\vec{a}^{\langle i \rangle}, b^{\langle i \rangle})$ tuples. Given this setup, the following hard problems constitute the defense mechanism of the LWE (Learning with Errors) cryptosystem:

$ $

\begin{itemize}
\item \textbf{Search problem:} Given the samples, recover the secret $\vec s$. The search-LWE assumption states that no efficient algorithm succeeds with non-negligible probability for the selected parameter family.
\item \textbf{Decision problem:} Distinguish LWE samples $(\vec a,\langle\vec a,\vec s\rangle+e)$ from independent uniform pairs in $\mathbb Z_q^k\times\mathbb Z_q$. The decisional-LWE assumption states that every efficient distinguisher has negligible advantage for the selected parameter family.
\end{itemize}

$ $

Search-to-decision reductions require hypotheses; the two formulations are not unconditionally interchangeable for every modulus and distribution. For example, Regev's reduction in Lemma 4.2 assumes a prime modulus bounded by a polynomial in the dimension, with its stated error and oracle conditions. His worst-case lattice reduction also has explicit noise and modulus requirements and is a quantum reduction. See \href{https://cims.nyu.edu/~regev/papers/qcrypto.pdf}{Regev, On Lattices, Learning with Errors, Random Linear Codes, and Cryptography}, Theorem 1.1 and Section 4.

$ $

\textbf{\underline{RLWE Problem}}

For the algebraic RLWE form, use $B=AS+E$ in $\mathbb Z_q[X]/(X^n+1)$, with polynomial representatives of degree less than $n$. Search asks for $S$; decision distinguishes these pairs from uniform ring pairs. This notation alone does not transfer every LWE security theorem: the ring, modulus, and error/secret distributions must satisfy the hypotheses of a ring-LWE result.

\end{tcolorbox}




\noindent\textbf{Small-secret scope:} A normal-form reduction can replace a uniform LWE secret by a secret drawn from the \emph{error distribution}. For a prime-power modulus, choose an invertible mask matrix $A_0$ using only its masks, with $\vec b_0=A_0\vec s+\vec e_0$. For a fresh row sample $(\vec a,b)$ set
\[
\vec a'=-\vec aA_0^{-1},\qquad b'=b-\vec aA_0^{-1}\vec b_0
=\vec a'\vec e_0+e.
\]
The new mask is uniform and the new secret has distribution $\chi^k$; uniform input pairs also map to uniform pairs. The preprocessing data are internal to this reduction. In the ring version the same identity requires an invertible initial ring mask. See \href{https://web.cs.ucla.edu/~sahai/work/web/2009\%20Publications/ACPS09.pdf}{Applebaum--Cash--Peikert--Sahai, Section 3.1, Lemma 2}. This argument does not turn Gaussian errors into a uniformly ternary secret. The later ternary-secret constructions therefore use a decisional assumption for their actual distribution, unless a separate distribution-matching reduction is supplied. Neither a dimension count nor the adjective ``small'' supplies that reduction.

\subsection{LWE Cryptosystem}
\label{subsec:lattice-scheme}

For the following correctness formulas, assume $t\mid q$, set $\Delta=q/t$, and let $[x]_q$ denote the centered integer representative of $x\bmod q$. Rounding is performed on this representative, and the decoded result is reduced modulo $t$. All bit-shift descriptions apply only when $\Delta$ is a power of two; in general scaling means multiplication or division by $\Delta$. The complete ciphertext is $(\vec a,b)$, not $b$ alone.

The LWE cryptosystem uses the following encryption formula: $b = \vec{s} \cdot \vec{a} + \Delta \cdot m + e$ (where $\vec{s}$ is a secret key, $\vec{a}$ is a publicly known random vector picked per encryption, $m$ is a plaintext, $e$ is small noise randomly picked per encryption from a normal distribution, and $b$ is a ciphertext). $\Delta$ is a scaling factor of the plaintext $M$ (shifting $m$ by $\text{log}_2\Delta$ bits to the left). Before encrypting the plaintext, we left-shift the plaintext several bits (i.e., $\text{log}_2\Delta$ bits) to secure sufficient space to store the error in the lower bits.

\begin{figure}[h!]
    \centering
  \includegraphics[width=0.8\linewidth]{figures/TFHE-fig1.pdf}
  \caption{An illustration of LWE's plaintext scaling and adding a noise: $\Delta \cdot m + e \in \mathbb{Z}_q$}
  \label{fig:scaling}
\end{figure}

\autoref{fig:scaling} visually illustrates the term $\Delta \cdot m + e$, where the plaintext $m$ left-shifted by $\text{log}_2\Delta$ bits and noised by the noise $e$. The actual encryption and decryption formulas are as follows:

\begin{tcolorbox}[title={\textbf{\tboxlabel{\ref*{subsec:lattice-scheme}} Lattice-based LWE Cryptosystem}}]
\begin{itemize}
    \item \textbf{\underline{Encryption}: } $b^{\langle i \rangle} = \vec{s} \cdot \vec{a}^{\langle i \rangle} + \Delta \cdot m^{\langle i \rangle} + e^{\langle i \rangle}$, where $b^{\langle i \rangle}$ and $\vec{a}^{\langle i \rangle}$ are publicly known also to the attacker, while $\vec{s}, m^{\langle i \rangle}, e^{\langle i \rangle}$ are unknown (only known by the secret key owner). 

    $ $
    
    \item \textbf{\underline{Decryption}: } $\left\lceil [b^{\langle i\rangle}-\vec s\cdot\vec a^{\langle i\rangle}]_q/\Delta\right\rfloor\bmod t=m^{\langle i\rangle}$ $\Big($ provided $|e^{\langle i \rangle}| < \dfrac{\Delta}{2}\Big)$
\end{itemize}
\end{tcolorbox}

$\lfloor \rceil_\Delta$ means rounding the number to the nearest multiple of $\Delta$. For example, $\lfloor 16 \rceil_{10} = 20$, which is rounding 16 to the nearest multiple of 10. As another example, $\lfloor 17 \rceil_{8} = 16$, which rounds 17 to the nearest multiple of 8 (note that 17 is closer to 16 than to 24; thus, it is rounded to 16). 


$ $

\noindent \textbf{Correctness: } In the decryption scheme, computing $b^{\langle i \rangle} - \vec{s} \cdot \vec{a}^{\langle i \rangle}$ gives $\Delta \cdot m^{\langle i \rangle}+ e^{\langle i \rangle}$, which is \autoref{fig:scaling}. Then, $\lceil \Delta \cdot m^{\langle i \rangle} + e^{\langle i \rangle} \rfloor_{\Delta}$ (i.e., rounding the value to the nearest multiple of $\Delta$) gives $\Delta \cdot m^{\langle i \rangle}$, provided the added noise $|e^{\langle i \rangle}| < \frac{\Delta}{2}$. That is, if the noise is less than $\frac{\Delta}{2}$, it will disappear during the rounding. Finally, right-shifting $\Delta \cdot m^{\langle i \rangle}$ by $\text{log}_2 \Delta$ bits gives $m^{\langle i \rangle}$. To summarize, if we ensure $|e^{\langle i \rangle}| < \frac{\Delta}{2}$ (which is why the noise $e^{\langle i \rangle}$ should be smaller than this threshold), then we can eliminate $e^{\langle i \rangle}$ during the decryption's rounding process and retrieve the original $\Delta \cdot m^{\langle i \rangle}$. The reason we scaled $m^{\langle i \rangle}$ by $\Delta$ is to: (i) create space for storing $e^{\langle i \rangle}$ in the lower bits during encryption such that the noise bits do not interfere with the plaintext bits (to avoid corrupting the plaintext bits); and (ii) blow away the noise $e^{\langle i \rangle}$ stored in the lower bits during decryption without corrupting the plaintext $m^{\langle i \rangle}$. 

$ $

\noindent \textbf{Security: } Fix the actual secret distribution, error distribution, modulus, dimension, and sample bound $Q$. Let $\delta(T,Q)$ be the maximum difference in a time-$T$ distinguisher's acceptance probabilities between a real LWE sample oracle and an independent uniform-pair oracle. The same secret is reused, while every query has a fresh independent uniform mask and fresh error. This is a decisional assumption, not merely a claim that key recovery is difficult.

In the adaptive left/right encryption game, a hidden bit $\vartheta$ selects one of each submitted pair $(m_0,m_1)$. Define the guessing bias $\epsilon=|\Pr[\widehat\vartheta=\vartheta]-1/2|$. A reduction chooses $\vartheta$ itself, takes $(\vec a,z)$ from its sample oracle for each query, and returns $(\vec a,z+\Delta m_\vartheta)$. With real samples this is exactly encryption. With uniform samples, the fresh $z$ remains uniform after any message offset, even conditional on the entire preceding transcript. Induction therefore makes the complete view independent of $\vartheta$. Outputting whether the adversary guessed correctly gives
\[
\epsilon\le\delta(T',Q),
\]
where $T'$ includes the simulation overhead. No extra factor $Q$ is needed because the assumed advantage already permits $Q$ samples. If IND advantage is instead defined as the output-probability difference between the two message experiments, the corresponding bound is $2\delta(T',Q)$. This covers repeated and adaptively chosen plaintexts through a classical encryption interface, without additional secret-dependent public material. For quantum attackers the underlying decisional assumption must cover quantum distinguishers; this statement does not grant superposition encryption queries.

Decryption subtracts the secret-dependent term, rounds to remove a sufficiently small error, divides by $\Delta$, and reduces modulo $t$. The noise bound and the security assumption are distinct requirements: small enough for correct decoding does not by itself mean large enough for security.

\subsection{RLWE Cryptosystem}
\label{subsec:lattice-scheme2}

In the RLWE cryptosystem, the formula in \tboxlabel{\ref*{subsec:lattice-scheme}} is the same, but $\vec{s}, \vec{a}^{\langle i \rangle}, b^{\langle i \rangle}, m^{\langle i \rangle}, e^{\langle i \rangle}$ are replaced by polynomials $S, A^{\langle i \rangle}, B^{\langle i \rangle}, M^{\langle i \rangle}, E^{\langle i \rangle}$ as follows:

\begin{tcolorbox}[title={\textbf{\tboxlabel{\ref*{subsec:lattice-scheme2}} Lattice-based RLWE Cryptosystem}}]
\begin{itemize}
    \item \textbf{\underline{Encryption}: } $B^{\langle i \rangle} = S \cdot A^{\langle i \rangle} + \Delta \cdot M^{\langle i \rangle} + E^{\langle i \rangle}$, where $B^{\langle i \rangle}$ and $A^{\langle i \rangle}$ are publicly known also to the attacker, while $S, M^{\langle i \rangle}, E^{\langle i \rangle}$ are unknown (only known by the secret key owner). 

    $ $
    
    \item \textbf{\underline{Decryption}: } $\left\lceil [B^{\langle i\rangle}-S\cdot A^{\langle i\rangle}]_q/\Delta\right\rfloor\bmod t=M^{\langle i\rangle}$ \\$\Big($provided $||E^{\langle i \rangle}||_{\infty} < \dfrac{\Delta}{2}$, meaning each coefficient of $E^{\langle i \rangle}$ has a magnitude less than $\dfrac{\Delta}{2}\Big)$

    $ $

    $\lceil  \rfloor_{\Delta}$ is equivalent to rounding each term's coefficient in the polynomial. 

\end{itemize}



\end{tcolorbox}

The security reduction just given applies to ring pairs under the corresponding decisional RLWE assumption, and to $k$ ring masks under the actual GLWE assumption. Expanding a ring product into coefficient equations gives a structured negacyclic matrix, not independent uniform LWE rows. Thus the proof template transfers, but a scalar LWE hardness theorem does not automatically certify a ring or module parameter family.
